\section{Online adaptation: protocol and resource measurements}
\label{app:adaptive-development}

\paragraph{Inputs and temporal roles.}
The eight synthetic streams contain 2,048 retained observations each, after
256 regime-A burn-in observations. Four seeds, 20261101--20261104, pair
stationary and recurring schedules by innovations. Histories and horizons
are respectively 12 and $(1,4,12)$. The retained Dominick series T7830,
T25851 and T108533 each contain 393 weekly observations, with histories 52
and horizons $(1,2,4)$. The Sunspot input is the finite final 4,096-point
crop of the previously selected daily T1 series, starting at original raw
index 69,828; histories are 28 and horizons $(1,3,7)$. No interpolation is
used. The source archives and version identifiers are documented with the
reproduction materials \citep{Godahewa2021Monash}.

At observation time $t$, the evaluator scores previously issued forecasts
whose targets mature at $t$, releases the newly available labels, updates
the model, then issues forecasts for $t+h$. Initial fitting uses common
joint-horizon complete rows, and calibration weights are selected before
processing the first evaluation observation using only earlier targets.
Matured earlier evaluation labels may update the model; the update rule and
selected weights remain fixed. The model therefore uses labels from earlier evaluation times once
they become available, while the update algorithm remains fixed. Scale factors use the fitting
prefix only. Losses average origins within horizon, horizons within series,
then series within collection; Dominick and Sunspot receive equal weight.
The four seeds define independent simulation replications; forecasts
from overlapping windows within a stream are dependent. No broad transfer intervals
are claimed from the two selected natural collections.

\paragraph{Reference policy and calibration scope.}
The geometric bank contains four local models of rank at most two. Every
16 observations, an adaptive arm considers one local model in round-robin
order, retains half its initial references as anchors, and replaces at most
four other references with observed histories assigned to that local model.
The affected geometry and all residual heads are rebuilt coherently from
their bounded replay. No bandwidth, penalty, history, reference-policy or
refresh search is performed. The kernel comparator uses an initial-prefix
median-distance bandwidth, regularizer 0.1 and FIFO supports, with recorded
inverse-residual checks and rebuilds. Calibration selects only the four
possible horizon-specific correction weights stated in
Section~\ref{sec:adaptive-development}.
For the kernel comparator, the correction is its kernel forecast minus
the common raw ridge forecast. The same calibration rule therefore selects
a ridge--KRR mixture for each horizon: weights one and zero recover KRR
and ridge, respectively. All reported method losses include this selection.

\paragraph{Prespecified accuracy and service criteria.}
The decision to expand the evaluation requires one capacity setting
to satisfy all accuracy and service criteria. On recurring
streams, adaptive KAHM must reduce mean loss by at least 3\% against each
of adaptive affine, adaptive RBF-local and kernel ridge, with wins on at
least three of four paired seeds against each, and by at least 1\% against
both fixed KAHM and raw ridge. On the natural aggregate, it must improve
each of the three main comparators by at least 1\%; at least one collection must show improvements of at least 3\%
against all three comparators, and neither collection may be more
than 5\% worse than any. KAHM must also meet the per-stream 99th-percentile service-time screen,
with no unexpected exception or cap violation in any method. These
thresholds were fixed before the original run and were also used for
the repeated evaluation after the numerical amendment, designated A1.
Neither capacity setting meets the recurring-stream accuracy criteria.
The protocol therefore classifies both settings as insufficient for
expanding the experiment. Every criterion is evaluated within a single
capacity setting.

% Generated from saved task-level resource measurements.
\begin{table}[!htbp]
\centering\small\setlength{\tabcolsep}{5pt}
\caption{Measured service and state by method/tier. P99 denotes the 99th percentile and is first computed per stream, then maximized across the 12 streams. Other columns also report maxima; they may occur on different streams. Owned NumPy buffers exclude Python objects and temporary work arrays; worker resident set size (RSS) includes interpreter, libraries and evidence storage. Neither is a precise incremental deployment footprint.}
\label{tab:adaptive-development-resources}
\begin{tabular}{@{}llrrrr@{}}
\toprule
Method & Tier & P99 (ms) & Max (ms) & NumPy (KiB) & RSS (MiB) \\
\midrule
Raw ridge & Small & 0.55 & 1.77 & 300.34 & 150.48 \\
Raw ridge & Large & 1.26 & 4.54 & 465.72 & 150.48 \\
Fixed KAHM & Small & 4.98 & 29.01 & 466.29 & 150.48 \\
Fixed KAHM & Large & 8.61 & 27.24 & 685.38 & 150.48 \\
Adaptive KAHM & Small & 9.89 & 49.10 & 518.95 & 150.48 \\
Adaptive KAHM & Large & 11.75 & 23.60 & 824.52 & 150.48 \\
Adaptive affine & Small & 9.56 & 24.75 & 509.83 & 150.48 \\
Adaptive affine & Large & 7.00 & 13.44 & 791.39 & 150.48 \\
Adaptive RBF-local & Small & 5.05 & 12.05 & 509.83 & 150.48 \\
Adaptive RBF-local & Large & 4.58 & 8.00 & 791.39 & 150.48 \\
Sliding-window KRR & Small & 3.15 & 3.95 & 477.34 & 150.48 \\
Sliding-window KRR & Large & 3.42 & 13.49 & 1011.72 & 150.48 \\
\bottomrule
\end{tabular}
\end{table}


\paragraph{Resource scope.}
Workers execute serially in fresh processes with one numerical thread and
the same CPU affinity. Initialization and calibration are measured
separately. The service interval covers history/pair assembly, all matured
horizon updates, feature/head rebuilds, kernel inverse audits and prediction;
scoring, serialization and state inspection are excluded. Buffer-state
samples and whole-worker peak RSS have the distinct scopes stated in
Table~\ref{tab:adaptive-development-resources}. The 20\,ms allowance is a
host-specific engineering screen, not an application-derived service-level
requirement. Different support sharing across horizons prevents interpreting
the tier labels as equal memory or equal computation.

\paragraph{Solver amendment and computational verification.}
The original implementation accumulated correction cross-products;
large rows entering and leaving a finite window produced cancellation in
the reproduced failing state. The amended implementation (A1) recomputes residuals against the
current baseline from the same bounded replay, centers to eliminate the
free intercept, equilibrates both data and penalty columns, and solves
the augmented ridge objective with pivoted QR. It changes arithmetic, not
the specified objective or advancement rule. The original 144 runs remain
in \path{extensions/adaptive_development/archives/initial_attempt.zip}.
Because the amended runs repeat the same inputs, they do not increase
the number of independent experimental replications. The verification program passes
\ADAuditChecks{} computational consistency checks, covering chronology,
coverage, calibration selection, aggregation, decision gates, resources,
shared baselines and matched reference changes. It imports neither the
forecasting model nor its analyzer, but remains an internal computational
check rather than independent scientific replication.

The reproduction materials contain both versions of the protocol and
source, the numerical diagnostic states, input hashes, predictions,
update traces, and summaries. The manuscript builder
\path{extensions/tmlr_integration/build_adaptive_assets.py} reconstructs
the displayed collection means and resource maxima from saved task
records, checks them against the analysis and decision files, and records
input/output hashes. It performs no model fitting. Both execution checkpoints were recorded locally. The study was not
publicly preregistered.
